\subsection{概形的下降}\label{src:3.1.3-subsection}
\begin{proposition}[开子概形与闭子概形的 fpqc 下降]\label{src:3.1.11}
设 $f:S'\to S$ 是概形的 fpqc 态射. 若 $Z'\subseteq S'$ 是闭 (或开) 子概形, 且在 $S'\times_S S'$ 中有子概形等式 $p_1^{-1}(Z')=p_2^{-1}(Z')$, 则存在闭 (或开) 子概形 $Z\subseteq S$, 使得 $Z'=f^{-1}(Z)$.
\end{proposition}
\begin{proof}
若 $Z'\hookrightarrow S'$ 是由理想层 $\mathcal I_{Z'}\subseteq\mathcal O_{S'}$ 定义的闭浸入, 则由拟凝聚层的 fpqc 下降~(\bookref{3.1.4}), $\mathcal I_{Z'}$ 下降为 $S$ 上的拟凝聚层 $\mathcal I_Z$, 且包含映射 $\mathcal I_{Z'}\hookrightarrow\mathcal O_{S'}$ 下降为包含映射 $\mathcal I_Z\hookrightarrow\mathcal O_S$. 因而 $Z'$ 下降为由 $\mathcal I_Z$ 定义的闭子概形 $Z\subseteq S$. 开浸入的情形可以通过取约化补集来处理.
\end{proof}
在下面的结果中, 用 $f^*X$ 表示沿态射 $f:S'\to S$ 对 $X\to S$ 作基变换, 会更加方便.
\begin{proposition}[仿射概形与拟仿射概形的 fpqc 下降]\label{src:3.1.12}
设 $f:S'\to S$ 是概形的 fpqc 态射. 若 $X'\to S'$ 是仿射 (或拟仿射) 态射, 且 $\alpha:p_1^*(X')\xrightarrow{\sim}p_2^*(X')$ 是 $S'\times_S S'$ 上满足 $p_{23}^*\alpha\circ p_{12}^*\alpha=p_{13}^*\alpha$ 的同构, 则存在概形的仿射 (或拟仿射) 态射 $X\to S$, 以及 $S'$ 上的同构 $\varphi:X'\to f^*(X)$, 使得 $p_1^*\varphi=p_2^*\varphi\circ\alpha$.
\end{proposition}
换言之, 存在虚线箭头将下图补全:
\[
\begin{tikzcd}[column sep=large,row sep=large]
{p_{23}^*\alpha\circ p_{12}^*\alpha=p_{13}^*\alpha}\arrow[d,dash]&
{p_1^*X'\xrightarrow{\alpha}p_2^*X'}\arrow[d]&
X'\arrow[d]\arrow[r,dashed]&X\arrow[d,dashed]\\
S'\times_S S'\times_S S'
\arrow[r,shift left=2.2ex,"p_{12}"]
\arrow[r,"p_{13}" description]\arrow[r,shift right=2.2ex,"p_{23}"']&
S'\times_S S'\arrow[r,shift left=1.1ex,"p_1"]
\arrow[r,shift right=1.1ex,"p_2"']&S'\arrow[r,"f"]&S
\end{tikzcd}
\]
\begin{proof}
若 $X'\to S'$ 仿射, 可将其写成 $X'=\Spec_X\mathcal A'$, 其中 $\mathcal A'$ 是拟凝聚的 $\mathcal O_{S'}$-代数层. 拟凝聚层的 fpqc 下降~(\bookref{3.1.4})~使我们可以先将 $\mathcal A'$ 下降为 $X$ 上的拟凝聚层 $\mathcal A$, 再将乘法映射 $\mathcal A'\otimes_{\mathcal O_{S'}}\mathcal A'\to\mathcal A'$ 下降为映射 $\mathcal A\otimes_{\mathcal O_S}\mathcal A\to\mathcal A$. 由下降, 这个映射必然满足使 $\mathcal A$ 成为拟凝聚 $\mathcal O_S$-代数的公理. 因而 $X'\to S'$ 下降为仿射态射 $X:=\Spec_S\mathcal A\to S$.

对于拟仿射态射, 利用典范分解 $X'\hookrightarrow\Spec_{S'}g_*'\mathcal O_{X'}\to S'$, 将其写成一个开浸入后接一个仿射态射, 再将上述仿射情形与开子概形的 fpqc 下降~(\bookref{3.1.11})~结合即可. 另见~\cite[定理 $6.6$]{BLR90}~与~\cite[定理 $4.33$]{Vis05}.
\end{proof}
\begin{remark}[上循环条件]\label{src:3.1.13}
上循环条件是概形下降的必要条件. Shimura 证明, 对一般的 $a,b\in\mathbb C$, 由
\[
y^2=x^6+ax^5+bx^4+x^3-bx^2+ax-1
\]
定义的 $\mathbb C$ 上的亏格 $2$ 曲线 $C'$ 同构于其复共轭, 却不能下降为实曲线~\cite{Shi72}. 换言之, 对于覆盖 $S'=\Spec\mathbb C\to\Spec\mathbb R$, 存在同构 $\alpha:p_1^*C'\xrightarrow{\sim}p_2^*C'$, 但 $C'\to S'$ 不能下降为曲线 $C\to S$.
\end{remark}
\begin{theorem}[局部拟有限且分离概形的 fppf 下降]\label{src:3.1.14}
设 $f:S'\to S$ 是概形的 fppf 态射. 若 $X'\to S'$ 是概形的局部拟有限且分离态射, 且 $\alpha:p_1^*(X')\xrightarrow{\sim}p_2^*(X')$ 是 $S'\times_S S'$ 上满足 $p_{23}^*\alpha\circ p_{12}^*\alpha=p_{13}^*\alpha$ 的同构, 则存在概形的局部拟有限且分离态射 $X\to S$, 以及 $S'$ 上的同构 $\varphi:X'\to f^*(X)$, 使得 $p_1^*\varphi=p_2^*\varphi\circ\alpha$.
\end{theorem}
这将表明, $\Sch$ 上参数化局部拟有限且分离态射 $X\to S$ 的预叠 $\mathsf{SepLQFin}$ 是 fppf 拓扑中的叠.
\begin{proof}
我们的策略是归约到拟仿射态射的 fpqc 下降~(\bookref{3.1.12}). 对每个拟紧开子集 $U'\subseteq X'$, 由 Zariski 主定理~(\bookref{A.7.3}), 复合 $U'\hookrightarrow X'\to S$ 是拟仿射的. 由于 $S'\to S$ 是 fppf 的, $p_2:S'\times_S S'\to S'$ 也是如此, 因而像 $V':=p_2(\alpha(p_1^*U'))$ 是 $X'$ 的开子集. 像 $V'$ 也是拟紧的, 包含 $U'$, 而且 $\alpha$ 限制为满足上循环条件的同构 $p_1^*(V')\xrightarrow{\sim}p_2^*(V')$. 映射 $V'\to S'$ 是拟仿射的, 由拟仿射态射的 fpqc 下降, $V'\to S'$ 下降为拟仿射态射 $V\to S$.

用拟紧开子集 $U_i'$ 覆盖 $X'$, 则子集 $V_i':=p_2(\alpha(p_1^*U_i'))$ 也覆盖 $X'$. 每个 $V_i'$ 都下降为 $S$ 上的拟仿射概形 $V_i$, 因而可以将概形 $V_i$ 粘合为 $S$ 上的概形 $X$, 其拉回即为 $X'$. 由直接的下降论证可知, $X\to S$ 局部拟有限且分离; 见命题~\bookref{3.1.26}. 另见~\cite[引理 $5.7.2$]{GR71}~与~\cite[Tag 02W7]{SP}.
\end{proof}
\begin{remark}[历史说明]\label{src:3.1.15}
\begingroup
\setlength{\emergencystretch}{1em}
若 $X'\to S'$ 拟紧, 则它是拟仿射的. 此时有效下降的论证是直接的 (见命题~\bookref{3.1.12}). 在 $1963$ 年给 Mumford 的一封信中~\cite[第 681 页]{Mum10}, Grothendieck 表示他能够证明局部拟有限且分离态射的有效下降. 这一结果出现于 Murre 未发表的工作中, 后来在局部 Noether 假设下发表于~\cite[命题 $1$]{Mur95}, 而在局部有限表现假设下出现于~\cite[引理 X.5.4]{SGA3II}. Raynaud 与 Gruson 在完全一般的情形下证明了这一定理, 但遗憾的是, 它被隐藏在~\cite[引理 $5.7.2$]{GR71}~的证明中. 当时这一定理没有广为人知, 这或许也是~\cite{Knu71,Art74b}~最初定义代数空间与代数叠时对对角态射要求拟紧的原因之一.

\par\endgroup
我们将经常以命题~\bookref{3.3.17}~的形式应用上述结果, 证明大平展或 fppf 拓扑中的一个给定层可由概形表出.
\end{remark}
\begin{example}[非有效下降]\label{src:3.1.16}
概形之间的任意态射 $X'\to S'$, 即使带有沿 fppf (甚至平展) 态射 $S'\to S$ 的下降数据, 也未必能下降为概形态射 $X\to S$. Hironaka 关于一个固有但非射影的三维簇的例子可以改造为一个不能下降的射影态射 $X'\to S'$; 见~\cite[Tag 08KE]{SP}. 此外, Raynaud 构造了一个维数为 $2$ 的正规 Noether 局部环 $A$, 一个平展覆盖 $S'\to S=\Spec A$, 以及一个不能下降为族 $C\to S$ 的光滑亏格 $1$ 曲线族 $C'\to S'$~\cite[XIII 3.2]{Ray70}. 还有如下例子: 一个离散赋值环 $R$, 一个平展覆盖 $S'\to S=\Spec R$, 以及一个一般纤维光滑而不能下降为族 $C\to S$ 的结点亏格 $1$ 曲线族 $C'\to S'$~\cite[§6.7]{BLR90}; 以及一个射影曲面 $S$, 一个平展覆盖 $S'\to S$, 以及一个一般纤维光滑而不能下降为族 $C\to S$ 的结点亏格 $0$ 曲线族 $C'\to S'$~\cite[例 $2.3$]{Ful10}.

另一方面, 带有沿 fppf 覆盖 $S'\to S$ 的下降数据的概形映射 $X'\to S'$ 总能下降为代数空间态射 $X\to S$ (平展情形见定理~\bookref{4.4.13}, 一般情形见推论~\bookref{7.3.6}). 事实上, 预叠 $\mathsf{AlgSp}$ 是 fppf 拓扑中的叠, 其中概形 $S$ 上的对象是代数空间态射 $X\to S$ (见习题~\bookref{5.5.15}).

不过, 有效下降在其他一些情形下也成立. 例如, 下一个结果将表明它对主 $G$-丛成立. 它也对二元组 $(X'\to S',\mathcal L')$ 成立, 其中 $X'\to S'$ 是概形的拟紧态射, 而 $\mathcal L'$ 是 $X'$ 上相对于 $S'$ 丰沛的线丛; 见~\cite[定理 $6.7$]{BLR90}~与~\cite[定理 $4.38$]{Vis05}.
\end{example}
\begin{proposition}[主 $G$-丛的 fpqc 下降]\label{src:3.1.17}
设 $G\to T$ 是 fppf 仿射群概形, $f:S'\to S$ 是 $T$ 上概形的 fpqc 态射. 若 $P'\to S'$ 是主 $G$-丛, 且 $\alpha:p_1^*P'\xrightarrow{\sim}p_2^*P'$ 是 $S'\times_S S'$ 上满足 $p_{12}^*\alpha\circ p_{23}^*\alpha=p_{13}^*\alpha$ 的主 $G$-丛同构, 则存在主 $G$-丛 $P\to S$, 以及主 $G$-丛同构 $\varphi:P'\to f^*P$, 使得 $p_1^*\varphi=p_2^*\varphi\circ\alpha$.
\end{proposition}
\begin{proof}
由于 $G\to T$ 仿射, 主 $G$-丛 $P'\to S'$ 仿射. 由仿射概形的 fpqc 下降 (命题~\bookref{3.1.12}), 存在仿射态射 $P\to S$ 及概形同构 $\varphi:P\to f^*Q$, 使得 $p_1^*\varphi=p_2^*\varphi\circ\alpha$. 由态射的 fpqc 下降~(\bookref{3.1.7}), 可以将作用 $G\times_S P'\to P'$ 下降为作用 $G\times_S P\to P$, 从而赋予 $P$ 主 $G$-丛结构, 并使 $\varphi:P'\to f^*P$ 成为 $G$-等变同构.
\end{proof}
